\section{Slutsats}\label{sec:conclusions}
Projektets mål var att utveckla och utvärdera hur kunskapsgrafer kan hjälpa en AI-agent att navigera och hämta intern företagsinformation, samt hur detta presterar gentemot mer traditionella system bestående av vanliga databaser och data lagrad i olika interna system.


\subsection{Framtida arbete}\label{subsec:future_work}
You should also explain potential future work based on your work.

% - Intressant att testa en lösning där API-baserae själv resonerar och anropar olika API:er med egen skapade API-calls
% - Det går att fixa API-baserade lösningen bättre genom att manuellt sitta och skapa varenda tool. Mer tools = Större risk att träffa rätt
% - Resonera kring hur man bygger en lösning som manuellt skapar KG genom API-anrop.
% - Finns en lösning som är mer hybrid? Te.x först kollar kunskapsgraf, och den inte hittar kollar den api:er? Om den hittar kan den sen bygga rätt koppling själv i grafen?



\subsection{Implementation av AI baserad KG'skapande/uppdaterande}\label{subsec:example_future_work}
Remember to not just point out potential future work, also explain how you would
approach this particular future work if you had the opportunity to do so.


